% GENERATED FILE. Edit canonical lessons, then run npm run generate:books.
% canonical-source-hash: fnv1a64:8465fef2d8c92365
% canonical-lessons: JA-C404-fuukei, JA-C404-nagameru, JA-C404-miwatasu, JA-C404-subarashii, JA-C404-miorosu

\chapter{The View from the Top}
\label{ch:ja-the-view-from-the-top}
\hlchaptermodality{404}

\begin{chapteropening}
\textbf{By the end of this chapter:} \emph{I can describe a view.}

Complete the last lesson of chapter 404: I can describe a view.
\end{chapteropening}

\section[\texorpdfstring{\ja{ふうけい} (fūkei)}{ふうけい (fūkei)}]{\ja{ふうけい} (fūkei) — scenery; a landscape}

\label{lesson:JA-C404-fuukei}



\begin{quote}
Before the new one: say the Japanese for a place visited on a trip, then the Japanese for to visit a place.
\end{quote}

\subsection*{You'll want to know: \ja{ふうけい}}
\textbf{\ja{ふうけい}} — \emph{fūkei} — \textquotedblleft{}scenery; a landscape\textquotedblright{}.

\begin{quote}
\textbf{\ja{いなかの} \ja{ふうけいは} \ja{うつくしいです。}} — \emph{inaka no fūkei wa utsukushii desu} — \textquotedblleft{}The countryside scenery is beautiful.\textquotedblright{}
\end{quote}

\begin{cousinweb}
\textbf{\ja{ふうけい}} (\emph{fūkei}) is Sino-Japanese (kango): two Chinese characters read with their on'yomi, readings borrowed from Chinese. They are \emph{fū}, \textquotedblleft{}style\textquotedblright{} (Middle Chinese \emph{pjuwng}) and \emph{kei}, \textquotedblleft{}scenery\textquotedblright{} (\emph{kjaengX}) (KANJIDIC2; Middle Chinese in Baxter's transcription).
\end{cousinweb}

\begin{grammarlens}[title={What you've built}]
One more word for getting around.
\end{grammarlens}

\subsection*{Your turn}
\begin{itemize}
\raggedright
  \item \emph{Say it:} \emph{fūkei}
  \item \emph{Say it:} \emph{fūkei}, once more
  \item \emph{Say it:} \emph{inaka no fūkei wa utsukushii desu} — The countryside scenery is beautiful.
  \item \emph{Recall:} say the Japanese for a place visited on a trip, then the Japanese for to visit a place, then say \emph{fūkei} again
\end{itemize}

\begin{tcolorbox}[breakable,colback=teal!4,colframe=teal!35!black,title={Before you move on}]
What does \ja{ふうけい} mean? (\textquotedblleft{}Scenery; a landscape\textquotedblright{}.) Say it once more.
\end{tcolorbox}
\section[\texorpdfstring{\ja{ながめる} (nagameru)}{ながめる (nagameru)}]{\ja{ながめる} (nagameru) — to gaze at; to admire (a view)}

\label{lesson:JA-C404-nagameru}



\begin{quote}
Before the new one: say the Japanese for to visit a place, then the Japanese for scenery.
\end{quote}

\subsection*{You'll want to know: \ja{ながめる}}
\textbf{\ja{ながめる}} — \emph{nagameru} — \textquotedblleft{}to gaze at; to admire (a view)\textquotedblright{}.

\begin{quote}
\textbf{\ja{まどから} \ja{うみを} \ja{ながめました。}} — \emph{mado kara umi o nagamemashita} — \textquotedblleft{}I gazed at the sea from the window.\textquotedblright{}
\end{quote}

Polite \textbf{\ja{ながめます}}, \emph{nagamemasu}; an ichidan verb. Past (polite): \textbf{\ja{ちちが} \ja{うみを} \ja{ながめました。}}, \emph{chichi ga umi o nagamemashita}.

\begin{cousinweb}
\textbf{\ja{ながめる}} (\emph{nagameru}) is a native Japanese word (wago), inherited rather than borrowed from Chinese; its kanji was attached later as a kun'yomi (JMdict).
\end{cousinweb}

\begin{grammarlens}[title={What you've built}]
One more word for getting around.
\end{grammarlens}

\subsection*{Your turn}
\begin{itemize}
\raggedright
  \item \emph{Say it:} \emph{nagameru}
  \item \emph{Say it:} \emph{nagameru}, once more
  \item \emph{Say it:} \emph{mado kara umi o nagamemashita} — I gazed at the sea from the window.
  \item \emph{Recall:} say the Japanese for to visit a place, then the Japanese for scenery, then say \emph{nagameru} again
\end{itemize}

\begin{tcolorbox}[breakable,colback=teal!4,colframe=teal!35!black,title={Before you move on}]
What does \ja{ながめる} mean? (\textquotedblleft{}To gaze at; to admire (a view)\textquotedblright{}.) Say it once more.
\end{tcolorbox}
\section[\texorpdfstring{\ja{みわたす} (miwatasu)}{みわたす (miwatasu)}]{\ja{みわたす} (miwatasu) — to look out over; to survey}

\label{lesson:JA-C404-miwatasu}



\begin{quote}
Before the new one: say the Japanese for scenery, then the Japanese for to gaze at.
\end{quote}

\subsection*{You'll want to know: \ja{みわたす}}
\textbf{\ja{みわたす}} — \emph{miwatasu} — \textquotedblleft{}to look out over; to survey\textquotedblright{}.

\begin{quote}
\textbf{\ja{やまの} \ja{うえから} \ja{まちを} \ja{みわたしました。}} — \emph{yama no ue kara machi o miwatashimashita} — \textquotedblleft{}From the top of the mountain I looked out over the town.\textquotedblright{}
\end{quote}

Polite \textbf{\ja{みわたします}}, \emph{miwatashimasu}; a godan verb. Past (polite): \textbf{\ja{あにが} \ja{まちを} \ja{みわたしました。}}, \emph{ani ga machi o miwatashimashita}.

\begin{cousinweb}
\textbf{\ja{みわたす}} (\emph{miwatasu}) is a native Japanese word (wago), inherited rather than borrowed from Chinese, built from \emph{mi}, \textquotedblleft{}seeing\textquotedblright{} (from \emph{miru}) and \emph{watasu}, \textquotedblleft{}to carry across\textquotedblright{}; its kanji were attached later as kun'yomi (JMdict).
\end{cousinweb}

\begin{grammarlens}[title={What you've built}]
One more word for getting around.
\end{grammarlens}

\subsection*{Your turn}
\begin{itemize}
\raggedright
  \item \emph{Say it:} \emph{miwatasu}
  \item \emph{Say it:} \emph{miwatasu}, once more
  \item \emph{Say it:} \emph{yama no ue kara machi o miwatashimashita} — From the top of the mountain I looked out over the town.
  \item \emph{Recall:} say the Japanese for scenery, then the Japanese for to gaze at, then say \emph{miwatasu} again
\end{itemize}

\begin{tcolorbox}[breakable,colback=teal!4,colframe=teal!35!black,title={Before you move on}]
What does \ja{みわたす} mean? (\textquotedblleft{}To look out over; to survey\textquotedblright{}.) Say it once more.
\end{tcolorbox}
\section[\texorpdfstring{\ja{すばらしい} (subarashii)}{すばらしい (subarashii)}]{\ja{すばらしい} (subarashii) — wonderful; splendid}

\label{lesson:JA-C404-subarashii}



\begin{quote}
Before the new one: say the Japanese for to gaze at, then the Japanese for to look out over.
\end{quote}

\subsection*{You'll want to know: \ja{すばらしい}}
\textbf{\ja{すばらしい}} — \emph{subarashii} — \textquotedblleft{}wonderful; splendid\textquotedblright{}.

\begin{quote}
\textbf{\ja{やまの} \ja{うえの} \ja{ふうけいは} \ja{すばらしかったです。}} — \emph{yama no ue no fūkei wa subarashikatta desu} — \textquotedblleft{}The scenery from the mountaintop was wonderful.\textquotedblright{}
\end{quote}

An i-adjective.

\begin{cousinweb}
\textbf{\ja{すばらしい}} (\emph{subarashii}) is a native Japanese word (wago), inherited rather than borrowed from Chinese; its kanji was attached later as a kun'yomi (JMdict).
\end{cousinweb}

\begin{grammarlens}[title={What you've built}]
One more word for getting around.
\end{grammarlens}

\subsection*{Your turn}
\begin{itemize}
\raggedright
  \item \emph{Say it:} \emph{subarashii}
  \item \emph{Say it:} \emph{subarashii}, once more
  \item \emph{Say it:} \emph{yama no ue no fūkei wa subarashikatta desu} — The scenery from the mountaintop was wonderful.
  \item \emph{Recall:} say the Japanese for to gaze at, then the Japanese for to look out over, then say \emph{subarashii} again
\end{itemize}

\begin{tcolorbox}[breakable,colback=teal!4,colframe=teal!35!black,title={Before you move on}]
What does \ja{すばらしい} mean? (\textquotedblleft{}Wonderful; splendid\textquotedblright{}.) Say it once more.
\end{tcolorbox}
\section[\texorpdfstring{\ja{みおろす} (miorosu)}{みおろす (miorosu)}]{\ja{みおろす} (miorosu) — to look down on (from above)}

\label{lesson:JA-C404-miorosu}



\begin{quote}
Before the new one: say the Japanese for to look out over, then the Japanese for wonderful.
\end{quote}

\subsection*{You'll want to know: \ja{みおろす}}
\textbf{\ja{みおろす}} — \emph{miorosu} — \textquotedblleft{}to look down on (from above)\textquotedblright{}.

\begin{quote}
\textbf{\ja{おかの} \ja{うえから} \ja{みなとを} \ja{みおろしました。}} — \emph{oka no ue kara minato o mioroshimashita} — \textquotedblleft{}From the top of the hill we looked down on the harbour.\textquotedblright{}
\end{quote}

Polite \textbf{\ja{みおろします}}, \emph{mioroshimasu}; a godan verb. Past (polite): \textbf{\ja{こどもが} \ja{おかから} \ja{まちを} \ja{みおろしました。}}, \emph{kodomo ga oka kara machi o mioroshimashita}.

\begin{cousinweb}
\textbf{\ja{みおろす}} (\emph{miorosu}) is a native Japanese word (wago), inherited rather than borrowed from Chinese, built from \emph{mi}, \textquotedblleft{}seeing\textquotedblright{} (from \emph{miru}) and \emph{orosu}, \textquotedblleft{}to lower\textquotedblright{}; its kanji were attached later as kun'yomi (JMdict).
\end{cousinweb}

\begin{grammarlens}[title={What you've built}]
One more word for getting around.
\end{grammarlens}

\subsection*{Your turn}
\begin{itemize}
\raggedright
  \item \emph{Say it:} \emph{miorosu}
  \item \emph{Say it:} \emph{miorosu}, once more
  \item \emph{Say it:} \emph{oka no ue kara minato o mioroshimashita} — From the top of the hill we looked down on the harbour.
  \item \emph{Recall:} say the Japanese for to look out over, then the Japanese for wonderful, then say \emph{miorosu} again
\end{itemize}

\begin{tcolorbox}[breakable,colback=teal!4,colframe=teal!35!black,title={Before you move on}]
What does \ja{みおろす} mean? (\textquotedblleft{}To look down on (from above)\textquotedblright{}.) Say it once more.
\end{tcolorbox}
